\documentclass[11pt,a4paper]{article}

\usepackage[margin=24mm]{geometry}
\usepackage[T1]{fontenc}
\usepackage{amsmath}
\usepackage{booktabs}
\usepackage{listings}
\usepackage[hidelinks]{hyperref}
\lstset{
    language=Python,
    basicstyle=\ttfamily\footnotesize,
    breaklines=true,
    columns=fullflexible,
    keepspaces=true,
    showstringspaces=false,
    frame=single
}
\setlength{\parskip}{4pt}
\setlength{\parindent}{0pt}

\title{Labwork 4 Report: Threads and 2D CUDA Blocks}
\author{Dao Quy Tung \\ Student ID: 2540109}
\date{3 October 2026}

\begin{document}
\maketitle

\section{Objective}
For this lab, I started with the grayscale code from Labwork 3 and changed
the GPU kernel to use 2D blocks.
I then used \texttt{time.time()} to compare the CPU and GPU running times
and calculate the speedup.

\section{Implementation}
I used the same $1200\times600$ RGB image as in Labwork 3,
\texttt{2021-Road-Race-Grand-Championship.jpg}. The grayscale calculation
also stays the same: take the average of the three colour values and
round down to an integer.
\[
\mathrm{gray}=\left\lfloor\frac{R+G+B}{3}\right\rfloor.
\]
On the CPU, a loop processes the pixels one by one after the image is
flattened. In both versions, I convert the channel values to integers
before adding them to avoid 8-bit overflow. The calculated grayscale
value is then written to all three channels.

In the GPU version, each thread handles one pixel. Its column and row
are given by \texttt{tidx} and \texttt{tidy}, so the image is accessed as
\texttt{src[tidy, tidx, channel]}. The \texttt{if} statement checks that
the pixel is inside the image before reading or writing it.
\begin{lstlisting}
@cuda.jit
def grayscale(src, dst):
    tidx = cuda.threadIdx.x + cuda.blockIdx.x * cuda.blockDim.x
    tidy = cuda.threadIdx.y + cuda.blockIdx.y * cuda.blockDim.y
    if tidy < src.shape[0] and tidx < src.shape[1]:
        r = int(src[tidy, tidx, 0])
        g = int(src[tidy, tidx, 1])
        b = int(src[tidy, tidx, 2])
        gray = np.uint8((r + g + b) // 3)
        dst[tidy, tidx, 0] = gray
        dst[tidy, tidx, 1] = gray
        dst[tidy, tidx, 2] = gray
\end{lstlisting}

I used a $16\times16$ block, or 256 threads per block. The grid dimensions
are rounded up so that the threads cover the whole image. For this input,
the grid is $75\times38$:
\begin{lstlisting}
blockSize = (16, 16)
gridSize = (
    (imageWidth + blockSize[0] - 1) // blockSize[0],
    (imageHeight + blockSize[1] - 1) // blockSize[1],
)
\end{lstlisting}

\section{Timing procedure}
I measured CPU time around the grayscale loop. For the GPU, I first ran
the kernel once and waited for it to finish with
\texttt{cuda.synchronize()}. This warm-up keeps the initial compilation
out of the measured time. I then timed the following part of the program:
\begin{lstlisting}
start = time.time()
devInput = cuda.to_device(image)
devOutput = cuda.device_array(image.shape, np.uint8)
grayscale[gridSize, blockSize](devInput, devOutput)
hostOutput = devOutput.copy_to_host()
gpuTime = time.time() - start
\end{lstlisting}
This measurement includes transferring the input to the GPU, allocating
the output, running the kernel, and bringing the result back to the CPU.
The call to \texttt{copy\_to\_host()} waits for the result, so the GPU work
has finished when the timer stops. Loading the image, saving the PNG
files, and comparing the outputs are done outside the timed sections.
I calculated speedup by dividing CPU time by GPU time:
\[
S=\frac{T_{\mathrm{CPU}}}{T_{\mathrm{GPU}}}.
\]

\section{Measured results}
I ran \texttt{Labwork4.py} on an NVIDIA GeForce RTX 2050 on
3 October 2026. The results from this run were:
\begin{center}
\begin{tabular}{lr}
\toprule
Measurement & Result \\
\midrule
CPU time & 0.856109 s \\
GPU time, including data transfer & 0.002371 s \\
Speedup & $361.10\times$ \\
\texttt{np.array\_equal(hostCPU, hostOutput)} & \texttt{True} \\
\bottomrule
\end{tabular}
\end{center}

The GPU was faster in this run, with a speedup of $361.10\times$,
including the time needed to transfer the data. I also compared the
outputs using \texttt{np.array\_equal}, which returned \texttt{True}.
This confirms that all 720,000 pixels match in all three channels.
The two images are saved as \texttt{gray\_cpu.png} and
\texttt{gray\_gpu.png}.

The speedup is calculated from the full timing values before they are
rounded for display. These numbers come from a single run with
$16\times16$ blocks, so the measured times may change on another run.

\end{document}
